\begin{frame}{Índices de linha e coluna}
\begin{itemize}
\item Em \texttt{int dados[2][3]}, há duas linhas e três colunas.
\item Linhas válidas: 0 e 1; colunas válidas: 0, 1 e 2.
\item Total de elementos: $2 \times 3 = 6$.
\item Cada linha é um vetor; em C, as linhas ficam em sequência na memória.
\end{itemize}
\begin{center}
\begin{tabular}{cccccc}
\toprule
\texttt{[0][0]} & \texttt{[0][1]} & \texttt{[0][2]} & \texttt{[1][0]} & \texttt{[1][1]} & \texttt{[1][2]} \\
24 & 25 & 26 & 28 & 29 & 30 \\
\bottomrule
\end{tabular}
\end{center}
\begin{alertblock}{Ambos os limites precisam ser respeitados}
\texttt{dados[0][3]} é inválido, mesmo existindo uma linha seguinte. Não tente atravessar uma linha usando um índice de coluna fora do intervalo.
\end{alertblock}
\end{frame}
